\documentclass[letterpaper,11pt]{article}
\usepackage[margin=1in]{geometry}
\usepackage{times}
\usepackage{graphicx}
\usepackage{amsmath,amssymb,amsfonts}
\usepackage{booktabs}
\usepackage{multirow}
\usepackage{xcolor}
\usepackage[hidelinks]{hyperref}
\usepackage{natbib}
\usepackage{array}
\usepackage{adjustbox}

\title{ContamCal: Mapping the Possibility Frontier of\\Contamination-Aware Benchmark Calibration}

\author{
Anonymous Authors
}

\begin{document}
\maketitle

\begin{abstract}
Benchmark contamination---the leakage of test items into training data---inflates reported model performance by up to 11.8 percentage points, yet no existing method provides calibrated scores with principled uncertainty quantification.
We introduce ContamCal, a framework that augments Item Response Theory (IRT) with contamination covariates, fusing four detection signals via Reciprocal Rank Fusion into a contamination-augmented 4PL IRT model with Bayesian credible intervals.
Through a systematic identifiability analysis grounded in Fisher Information theory, we establish a \emph{possibility frontier} for contamination calibration: practical identifiability requires detection signal AUC $\geq 0.72$, a threshold that current signals meet only under heavy SFT-regime contamination of smaller models ($\leq$1.5B parameters).
A 144-condition simulation sweep confirms a sharp phase transition at AUC $\approx 0.70$, below which the contamination parameter $\gamma$ becomes confounded with model ability $\theta$.
Empirically, we demonstrate that (1) contamination detection signals degrade gradually---not abruptly---across training regimes, refuting clean regime boundaries; (2) signal quality is inversely correlated with model scale, creating a fundamental ``calibration paradox'' where the models most in need of calibration are those where calibration is least feasible; and (3) in the RL regime, all calibration methods (including DCR) produce errors comparable to or worse than uncalibrated scores.
These findings reframe the contamination calibration problem from a solvable estimation task to a problem with well-characterized impossibility regions, providing the community with an actionable identifiability map and concrete signal quality targets for future detection research.
\end{abstract}

\section{Introduction}

Large language model (LLM) benchmarks serve as the primary basis for model comparison and selection, yet their reliability is increasingly undermined by training data contamination~\citep{sainz2023nlp,balloccu2024leak}.
Recent work demonstrates that rephrased test items can inflate MMLU accuracy by 8--15 percentage points while evading standard n-gram decontamination~\citep{yang2023rethinking,oren2024proving}.
The problem is compounded by RL post-training (e.g., GRPO), which further transforms contamination signals and conceals evidence of data leakage.

Existing approaches address contamination through two disjoint strategies: \emph{detection} methods (Min-K\%++~\citep{shi2024detecting}, ConStat~\citep{dekoninck2024constat}, KDS~\citep{dong2024generalization}) flag potentially contaminated items but do not produce calibrated scores, while \emph{mitigation} strategies (option shuffling, rephrasing, dynamic benchmarks) attempt to eliminate contamination effects directly.
\citet{xu2024benchmark} systematically evaluate 20 mitigation strategies and find none reliably effective---motivating a fundamentally different approach.

We draw an analogy to causal inference, where the field progressed from attempting to eliminate confounders to statistically controlling for them.
ContamCal instantiates this shift for benchmark evaluation: rather than eliminating contamination, we \emph{measure} its effect and subtract it with proper uncertainty quantification.
The framework is grounded in educational measurement's item compromise IRT tradition~\citep{eckerly2017detecting,sinharay2017simultaneous} and explanatory IRT~\citep{deboeck2004explanatory}, with three structural adaptations for the LLM setting: (a) independent examinees $\to$ a single model with joint item distributions, (b) binary leakage $\to$ continuous exposure scores requiring measurement error propagation, and (c) stable ability $\to$ regime-conditional ability requiring regime-aware signal selection.

However, the central finding of this work is not that ContamCal \emph{solves} contamination calibration, but rather that it reveals the precise conditions under which calibration is and is not feasible.
Through Fisher Information analysis and a 144-condition simulation sweep, we establish a \emph{possibility frontier}: contamination calibration requires detection signal AUC $\geq 0.72$, a threshold that current detection technology meets only in narrow operating conditions (heavy SFT contamination of smaller models).
This is a \emph{negative but informative} result: we prove that under current signal quality, any post-hoc calibration method---including the recent DCR~\citep{yang2023dcr}---cannot reliably recover true model performance.

Our contributions are:
\begin{enumerate}
\item \textbf{Theoretical:} The first Fisher Information identifiability analysis for contamination calibration, yielding a practical AUC threshold of 0.72 and Cram\'er-Rao lower bounds for the contamination parameter $\gamma$ (Section~\ref{sec:theory}).
\item \textbf{Empirical:} A signal survival analysis across four training regimes (Base, SFT-Clean, SFT-Contam, RL) showing gradual, benchmark-dependent degradation that refutes clean regime boundaries (Section~\ref{sec:signal_survival}).
\item \textbf{Methodological:} The ContamCal framework itself, which provides meaningful calibration when signals are sufficient (SFT regime, larger benchmarks) and honestly reports high uncertainty when they are not (Section~\ref{sec:calibration}).
\end{enumerate}

\section{Related Work}

\paragraph{Contamination Detection.}
Min-K\%++~\citep{shi2024detecting} detects memorization via token-level perplexity patterns.
ConStat~\citep{dekoninck2024constat} compares performance on original vs.\ paraphrased items.
Self-Critique leverages models' meta-cognitive evaluation of their own knowledge sources.
KDS~\citep{dong2024generalization} and CDD~\citep{oren2024proving} provide additional detection signals.
These methods flag contamination but do not calibrate scores.

\paragraph{IRT for LLM Evaluation.}
\citet{lalor2024benchmarking} apply IRT to LLM benchmarking at AAAI 2024.
HiBayES~\citep{polo2024tinybenchmarks} provides Bayesian uncertainty for benchmark scores but ignores contamination.
ATLAS~\citep{sedoc2020item} and Growing Pains~\citep{vania2021comparing} use psychometric models for evaluation but without contamination covariates.

\paragraph{Contamination Mitigation and Calibration.}
\citet{xu2024benchmark} test 20 mitigation strategies and find none effective.
DCR~\citep{yang2023dcr} (EMNLP 2025) produces adjusted accuracy via fuzzy inference but lacks uncertainty bounds and RL handling.
LiveBench~\citep{white2024livebench} offers continuously refreshed benchmarks but requires perpetual curation.

\paragraph{Item Compromise IRT.}
In educational measurement, \citet{eckerly2017detecting} and \citet{sinharay2017simultaneous} model item compromise within IRT frameworks, providing the theoretical foundation for our contamination covariate approach.
Explanatory IRT~\citep{deboeck2004explanatory} extends standard IRT with person- and item-level covariates; our $\gamma$ parameter follows this tradition.

\section{Method: The ContamCal Framework}
\label{sec:method}

ContamCal operates in two stages: (1) contamination signal computation and fusion, and (2) contamination-augmented IRT estimation.

\subsection{Stage 1: Signal Pipeline and RRF Fusion}

We compute four contamination detection signals for each benchmark item:
\begin{itemize}
\item \textbf{Embedding Retrieval}: Semantic similarity between test items and training corpus passages, measured via cosine similarity of sentence embeddings.
\item \textbf{Min-K\%++}~\citep{shi2024detecting}: Token-level perplexity anomaly scores indicating memorization.
\item \textbf{ConStat}~\citep{dekoninck2024constat}: Performance gap between original and paraphrased item variants.
\item \textbf{Self-Critique}: Model self-evaluation of knowledge source provenance.
\end{itemize}

Signals are fused via Reciprocal Rank Fusion (RRF;~\citealp{cormack2009reciprocal}):
\begin{equation}
\text{RRF}(i) = \sum_{s=1}^{S} \frac{1}{k + r_s(i)}
\label{eq:rrf}
\end{equation}
where $r_s(i)$ is item $i$'s rank under signal $s$ and $k=60$ is the standard smoothing constant.
RRF is a \emph{design choice}, not a contribution: it avoids the meta-circularity of learned aggregators whose training data may itself be contaminated.

\subsection{Stage 2: Contamination-Augmented 4PL IRT}

We augment the four-parameter logistic (4PL) IRT model with a contamination covariate $\gamma$:
\begin{equation}
P(Y_{ij}=1 \mid \theta_j, \tau_i) = c_i + (d_i - c_i) \cdot \sigma\!\big(a_i(\theta_j + \gamma \cdot \tau_i - b_i)\big)
\label{eq:4pl_gamma}
\end{equation}
where $Y_{ij}$ is the binary response of model $j$ to item $i$, $\theta_j$ is model ability, $\tau_i$ is the latent exposure of item $i$, $a_i$ is discrimination, $b_i$ is difficulty, $c_i$ and $d_i$ are lower and upper asymptotes, and $\gamma$ is the contamination inflation parameter.
The calibrated score is obtained by setting $\gamma=0$ in the fitted model.

\paragraph{Per-benchmark $\gamma_k$ variant.}
We also fit a per-benchmark extension $\gamma_k$ allowing contamination effects to vary across benchmarks, motivated by the empirical observation that inflation varies up to $2.4\times$ across benchmarks (Section~\ref{sec:sft_injection}).

\paragraph{Measurement Error Model.}
The observed exposure score $x_i$ is a noisy proxy for latent $\tau_i$:
\begin{equation}
x_i \sim \mathcal{N}(\tau_i, \sigma^2_{\epsilon,i})
\end{equation}
where $\sigma^2_{\epsilon,i}$ is estimated from signal disagreement.
This propagates detection uncertainty into posterior credible intervals automatically.

\paragraph{Priors.}
$\theta_j \sim \mathcal{N}(0,1)$, $\gamma \sim \mathcal{N}(0,1)$, $a_i \sim \text{LogNormal}(0,0.5)$, $b_i \sim \mathcal{N}(0,1)$, $c_i \sim \text{Beta}(5,20)$, $d_i \sim \text{Beta}(20,5)$.

\section{Theoretical Analysis: The Possibility Frontier}
\label{sec:theory}

\subsection{Structural vs.\ Practical Identifiability}

The contamination-augmented IRT model (Eq.~\ref{eq:4pl_gamma}) is \emph{structurally} identifiable: given perfect exposure scores $\tau_i$, the parameter $\gamma$ is uniquely determined by the data.
However, \emph{practical} identifiability---the ability to estimate $\gamma$ with useful precision from finite, noisy data---depends critically on detection signal quality.

\subsection{Fisher Information Analysis}

We derive the Fisher Information matrix for $\gamma$ as a function of signal AUC, number of evaluated models $N$, and item count.
The Cram\'er-Rao lower bound provides a floor on the variance of any unbiased estimator of $\gamma$:

\begin{table}[t]
\centering
\caption{Cram\'er-Rao lower bounds for $\gamma$ by signal AUC. Practical identifiability requires $\text{CR bound} \leq 0.15$, corresponding to AUC $\geq 0.72$.}
\label{tab:cramer_rao}
\begin{tabular}{@{}lccc@{}}
\toprule
AUC & CR Bound (sd) & Identifiable? & Min $N$ for $\gamma$ sig. \\
\midrule
0.55 & 0.45 & No & $>$500 \\
0.60 & 0.35 & No & $>$300 \\
0.65 & 0.28 & No & $>$120 \\
0.70 & 0.15 & Marginal & $\sim$60 \\
0.80 & 0.09 & Yes & $\sim$30 \\
0.90 & 0.05 & Yes & $\sim$15 \\
\bottomrule
\end{tabular}
\end{table}

Table~\ref{tab:cramer_rao} reveals a practical identifiability threshold at AUC $\approx 0.72$: below this, the theoretical minimum standard deviation of $\hat{\gamma}$ exceeds 0.15, making contamination effects indistinguishable from estimation noise.

\subsection{Simulation Sweep}
\label{sec:sim_sweep}

We validate the Fisher Information bounds via a 144-condition simulation sweep (6 AUC levels $\times$ 4 $N_{\text{models}}$ $\times$ 3 priors $\times$ 2 exposure types, 10 seeds each).

\begin{table}[t]
\centering
\caption{Identifiability map: $\gamma$ recovery by AUC and $N_{\text{models}}$ (default prior $\mathcal{N}(0,1)$). A sharp phase transition occurs at AUC $\approx 0.70$.}
\label{tab:identifiability_map}
\begin{adjustbox}{max width=\columnwidth}
\begin{tabular}{@{}cccccc@{}}
\toprule
AUC & $N$ & $\gamma$ RMSE & $\theta$ corr & Coverage 95\% & CI Width (pp) \\
\midrule
0.55 & 50 & 0.52 & 0.48 & 62\% & 32.0 \\
0.60 & 50 & 0.42 & 0.58 & 68\% & 26.5 \\
0.65 & 50 & 0.31 & 0.70 & 74\% & 21.0 \\
0.70 & 50 & 0.19 & 0.82 & 82\% & 14.5 \\
0.80 & 50 & 0.10 & 0.92 & 91\% & 8.2 \\
0.90 & 50 & 0.06 & 0.96 & 94\% & 5.1 \\
\midrule
0.65 & 20 & 0.45 & 0.52 & 65\% & 28.5 \\
0.65 & 100 & 0.22 & 0.78 & 78\% & 16.8 \\
0.65 & 200 & 0.17 & 0.84 & 82\% & 13.5 \\
\bottomrule
\end{tabular}
\end{adjustbox}
\end{table}

Table~\ref{tab:identifiability_map} confirms the phase transition: $\gamma$ RMSE drops from 0.31 to 0.19 between AUC 0.65 and 0.70, coverage jumps from 74\% to 82\%, and $\theta$-$\gamma$ posterior correlation drops from 0.70 to 0.82 (ability recovery).
At AUC $\leq 0.65$, the model confounds ability with contamination ($\theta$-$\gamma$ corr $= 0.58$--$0.70$).

\subsection{Prior Sensitivity}

Below the identifiability threshold, results become prior-dominated:

\begin{table}[t]
\centering
\caption{Prior sensitivity by AUC level. At AUC $= 0.65$, the sensitivity index of $0.42$ indicates prior-dominated estimation; at AUC $= 0.80$, it drops to $0.08$ (data-dominated).}
\label{tab:prior_sensitivity}
\begin{tabular}{@{}lcccc@{}}
\toprule
AUC & $\mathcal{N}(0,0.5)$ MAE & $\mathcal{N}(0,1)$ MAE & $\mathcal{N}(0,2)$ MAE & Sensitivity \\
\midrule
0.55 & 12.5 & 10.8 & 9.5 & 0.58 \\
0.65 & 9.0 & 7.5 & 6.2 & 0.42 \\
0.70 & 5.2 & 4.8 & 4.5 & 0.15 \\
0.80 & 2.8 & 2.6 & 2.5 & 0.08 \\
\bottomrule
\end{tabular}
\end{table}

At AUC $= 0.65$, the prior sensitivity index is 0.42 (Table~\ref{tab:prior_sensitivity}): changing the prior from $\mathcal{N}(0,0.5)$ to $\mathcal{N}(0,2)$ shifts MAE by 2.8pp.
This means the Bayesian posterior at low AUC reflects the modeler's beliefs more than the data---a diagnostic we propose as a general-purpose test for under-identification.

\section{Experiments}

\subsection{Controlled SFT Injection}
\label{sec:sft_injection}

We inject benchmark items into Qwen2.5 models (0.5B, 1.5B, 3B) at five dosage levels (0\%, 5\%, 25\%, 50\%, 100\%) across four benchmarks (MMLU, GSM8K, ARC-Challenge, HumanEval), yielding 15 SFT training runs.

\begin{table}[t]
\centering
\caption{Score inflation (absolute accuracy) by model, benchmark, and dosage. Inflation is inversely correlated with model scale and varies $2.4\times$ across benchmarks at 100\% dosage.}
\label{tab:inflation}
\begin{adjustbox}{max width=\columnwidth}
\begin{tabular}{@{}llccccc@{}}
\toprule
Model & Benchmark & 0\% & 5\% & 25\% & 50\% & 100\% \\
\midrule
\multirow{4}{*}{Qwen-0.5B}
 & MMLU & 42.1 & 43.0 & 45.5 & 49.3 & 53.9 \\
 & GSM8K & 30.8 & 31.1 & 32.3 & 33.8 & 35.8 \\
 & ARC-C & 37.9 & 38.5 & 40.0 & 42.4 & 45.2 \\
 & HumanEval & 21.3 & 21.5 & 22.2 & 23.5 & 25.0 \\
\midrule
\multirow{4}{*}{Qwen-1.5B}
 & MMLU & 52.8 & 53.4 & 54.9 & 56.8 & 60.2 \\
 & GSM8K & 48.2 & 48.5 & 49.6 & 51.2 & 53.0 \\
 & ARC-C & 51.6 & 52.0 & 53.2 & 55.0 & 57.5 \\
 & HumanEval & 34.8 & 35.0 & 35.8 & 37.0 & 38.4 \\
\midrule
\multirow{4}{*}{Qwen-3B}
 & MMLU & 61.5 & 61.8 & 62.8 & 64.3 & 67.0 \\
 & GSM8K & 62.0 & 62.2 & 62.9 & 63.8 & 65.2 \\
 & ARC-C & 60.2 & 60.5 & 61.4 & 62.8 & 65.0 \\
 & HumanEval & 47.8 & 48.0 & 48.6 & 49.7 & 51.3 \\
\bottomrule
\end{tabular}
\end{adjustbox}
\end{table}

Table~\ref{tab:inflation} reveals two structural properties.
First, inflation is inversely correlated with model scale: 0.5B shows 11.8pp MMLU inflation at 100\% dosage versus 5.5pp for 3B, because smaller models have greater capacity for memorization relative to generalization.
Second, inflation varies $2.4\times$ across benchmarks (3.5pp for HumanEval vs.\ 8.5pp for MMLU at 100\% dosage on 0.5B), because reasoning and code benchmarks resist paraphrase memorization more than MCQ formats.
The 5\% dosage condition yields only 0.9pp inflation on MMLU, indistinguishable from seed noise (std $= 2.9$pp).
These non-uniform patterns motivate per-benchmark $\gamma_k$ and challenge the additive-$\gamma$ model specification.

\subsection{Signal Quality}
\label{sec:signals}

\begin{table}[t]
\centering
\caption{Signal AUC by dosage level (averaged across models and benchmarks). Only Embedding Retrieval at 100\% dosage approaches the identifiability threshold of 0.72.}
\label{tab:signal_auc_dosage}
\begin{tabular}{@{}lcccc@{}}
\toprule
Signal & 5\% & 25\% & 50\% & 100\% \\
\midrule
Embedding Retrieval & 0.53 & 0.61 & 0.67 & 0.71 \\
Min-K\%++ & 0.52 & 0.57 & 0.62 & 0.66 \\
ConStat & 0.51 & 0.55 & 0.60 & 0.63 \\
Self-Critique & 0.51 & 0.54 & 0.57 & 0.60 \\
\bottomrule
\end{tabular}
\end{table}

\begin{table}[t]
\centering
\caption{Signal AUC by model scale at 100\% dosage (MMLU). Detection quality degrades with scale because larger models' pretraining knowledge confounds contamination signals.}
\label{tab:signal_auc_scale}
\begin{tabular}{@{}lccc@{}}
\toprule
Signal & Qwen-0.5B & Qwen-1.5B & Qwen-3B \\
\midrule
Embedding Retrieval & 0.74 & 0.72 & 0.67 \\
Min-K\%++ & 0.70 & 0.67 & 0.62 \\
ConStat & 0.66 & 0.64 & 0.60 \\
Self-Critique & 0.55 & 0.61 & 0.64 \\
\bottomrule
\end{tabular}
\end{table}

Tables~\ref{tab:signal_auc_dosage}--\ref{tab:signal_auc_scale} establish a catch-22: high-dosage contamination is detectable (AUC 0.71 at 100\%) but has obvious calibration value, while low-dosage contamination needs detection but signals are indistinguishable from random (AUC 0.51--0.53 at 5\%).
All four signals at 5\% dosage are statistically indistinguishable from chance.
Signal quality degrades with model scale (Table~\ref{tab:signal_auc_scale}): Embedding AUC drops from 0.74 (0.5B) to 0.67 (3B), because larger models ``naturally know'' benchmark content, confounding contamination-specific signals.
Average pairwise signal correlation is 0.55, indicating high redundancy and limited complementarity.

\subsection{Signal Aggregation}
\label{sec:aggregation}

\begin{table}[t]
\centering
\caption{Aggregation method comparison at 100\% dosage. RRF fusion (0.68) underperforms the best single signal (Embedding, 0.71) due to high inter-signal correlation ($\bar{r}=0.55$).}
\label{tab:aggregation}
\begin{tabular}{@{}lcc@{}}
\toprule
Method & AUC & $\Delta$ vs.\ Best Single \\
\midrule
Embedding Retrieval (best single) & 0.71 & --- \\
Min-K\%++ & 0.66 & $-$0.05 \\
ConStat & 0.63 & $-$0.08 \\
Self-Critique & 0.60 & $-$0.11 \\
\midrule
RRF (4-signal) & 0.68 & $-$0.03 \\
Rank Average (4-signal) & 0.67 & $-$0.04 \\
Weighted RRF (calibrated) & 0.69 & $-$0.02 \\
2-signal degraded (CS+SC) & 0.62 & $-$0.09 \\
\bottomrule
\end{tabular}
\end{table}

\begin{table}[t]
\centering
\caption{RRF AUC across dosage levels. Fusion provides no benefit over the best single signal at any dosage.}
\label{tab:rrf_dosage}
\begin{tabular}{@{}lcccc@{}}
\toprule
Method & 5\% & 25\% & 50\% & 100\% \\
\midrule
RRF 4-signal & 0.53 & 0.58 & 0.63 & 0.68 \\
Best single (Embedding) & 0.53 & 0.61 & 0.67 & 0.71 \\
2-signal degraded & 0.51 & 0.55 & 0.57 & 0.62 \\
\bottomrule
\end{tabular}
\end{table}

RRF fusion provides no improvement over the best single signal (Table~\ref{tab:aggregation}): 4-signal RRF achieves AUC 0.68 versus Embedding's 0.71 ($\Delta = -0.03$).
The failure stems from high inter-signal correlation ($\bar{r} = 0.55$): the four signals share approximately 30\% of variance, leaving minimal information gain from fusion.
Even calibrated weighted RRF (0.69) does not significantly improve over Embedding alone ($p > 0.10$).
False positive injection degrades RRF AUC by 0.10 at 20\% FP rate, indicating fragility under realistic noise.
The 2-signal mode available for closed-weight models (ConStat + Self-Critique) yields only 0.62 AUC.

\subsection{Noise Propagation}
\label{sec:noise}

\begin{table}[t]
\centering
\caption{CI width and coverage under noise injection and measurement error model. The measurement error model achieves near-nominal coverage (92\%) at the cost of vacuously wide intervals (22pp).}
\label{tab:noise_prop}
\begin{tabular}{@{}lcccc@{}}
\toprule
Noise Level & $\gamma$ CI (pp) & $\theta$ CI (pp) & Cov.\ 95\% & Cov.\ 80\% \\
\midrule
Baseline (0x) & 9.5 & 7.2 & 0.76 & 0.58 \\
Matched (1x) & 14.8 & 11.5 & 0.84 & 0.68 \\
1.5x & 19.5 & 15.8 & 0.91 & 0.78 \\
2x & 26.0 & 21.5 & 0.96 & 0.88 \\
Meas.\ error model & 22.0 & 15.5 & 0.92 & 0.78 \\
\bottomrule
\end{tabular}
\end{table}

Table~\ref{tab:noise_prop} exposes a fundamental precision-coverage tradeoff.
The measurement error model ``fixes'' coverage from 0.76 to 0.92, but at the cost of 22pp CI width---a 60\% accuracy model's 95\% CI would span [49\%, 71\%], rendering model selection vacuous.
At 72\% of conditions, CI width exceeds 15pp; at 58\%, it exceeds 20pp.
This is not a model deficiency but an honest report: at current signal quality (Embedding FP $= 22\%$, FN $= 18\%$; Self-Critique FP $= 38\%$, FN $= 30\%$), the data simply cannot support precise calibration.
We formalize this as \emph{operational identifiability}: even when $\gamma$ is estimable, the resulting calibration may produce intervals too wide for actionable decisions.

\subsection{Core Calibration Results}
\label{sec:calibration}

\begin{table*}[t]
\centering
\caption{Calibration MAE (pp) by method and benchmark. 4PL+$\gamma$ does not significantly outperform linear regression ($p = 0.18$) or DCR ($p = 0.35$). Only vanilla 4PL and Emperor mitigation are significantly worse.}
\label{tab:calibration_main}
\begin{adjustbox}{max width=\textwidth}
\begin{tabular}{@{}lccccc@{}}
\toprule
Method & MMLU & GSM8K & ARC-C & HumanEval & Average \\
\midrule
4PL+$\gamma$ global (ours) & 5.2$\pm$0.8 & 8.5$\pm$1.5 & 6.2$\pm$1.0 & 7.3$\pm$1.8 & 6.8$\pm$1.1 \\
4PL+$\gamma_k$ per-benchmark & 5.0$\pm$0.9 & 8.2$\pm$1.6 & 5.9$\pm$1.1 & 7.0$\pm$2.2 & 6.5$\pm$1.2 \\
\midrule
Raw (no adjustment) & 8.5 & 7.8 & 8.2 & 9.5 & 8.5 \\
Remove-and-rescore & 6.5 & 7.8 & 6.8 & 7.7 & 7.2 \\
Single-signal 2PL+$\gamma$ & 5.5 & 7.5 & 5.8 & 7.2 & 6.5 \\
Vanilla 4PL (no $\gamma$) & 8.0 & 9.5 & 8.2 & 9.5 & 8.8 \\
Linear regression + CI & 5.0 & 7.2 & 5.5 & 7.1 & 6.2 \\
Emperor mitigation & 7.2 & 8.8 & 7.5 & 8.5 & 8.0 \\
DCR (EMNLP 2025) & 5.0 & 7.8 & 5.8 & 7.0 & 6.4 \\
\bottomrule
\end{tabular}
\end{adjustbox}
\end{table*}

\begin{table}[t]
\centering
\caption{Interval quality metrics for 4PL+$\gamma$ global. Coverage is systematically below nominal at all levels, indicating model misspecification or under-identified parameters.}
\label{tab:interval_quality}
\begin{adjustbox}{max width=\columnwidth}
\begin{tabular}{@{}lccccc@{}}
\toprule
Metric & MMLU & GSM8K & ARC-C & HumanEval & Avg \\
\midrule
Coverage 95\% & 0.80 & 0.72 & 0.78 & 0.74 & 0.76 \\
Coverage 80\% & 0.62 & 0.52 & 0.60 & 0.56 & 0.58 \\
Coverage 50\% & 0.38 & 0.30 & 0.36 & 0.34 & 0.35 \\
CI width 95\% (pp) & 8.2 & 11.5 & 9.0 & 9.2 & 9.5 \\
Kendall $\tau$ & 0.78 & 0.65 & 0.75 & 0.68 & 0.72 \\
\bottomrule
\end{tabular}
\end{adjustbox}
\end{table}

Table~\ref{tab:calibration_main} presents the central calibration results.
4PL+$\gamma$ achieves 6.8pp average MAE, reducing error by 1.7pp versus raw scores (8.5pp) but \emph{not} significantly outperforming linear regression (6.2pp, $p = 0.18$) or DCR (6.4pp, $p = 0.35$).
Only two of seven baselines are significantly worse: vanilla 4PL ($p = 0.02$) and Emperor mitigation ($p = 0.04$)---a low bar, as neither accounts for contamination.
DCR outperforms ContamCal on MMLU (5.0 vs.\ 5.2pp) and ARC-C (5.8 vs.\ 6.2pp) despite lacking uncertainty bounds.

The per-benchmark $\gamma_k$ variant yields a negligible 0.3pp improvement (6.5 vs.\ 6.8pp), with HumanEval failing to converge (ESS $= 85$, $\hat{R} = 1.12$) due to insufficient item count ($N = 164$).
The $\gamma_k$ posteriors overlap 78\% with the global $\gamma$ posterior, and a Wald test finds no significant difference.

\paragraph{Dose-response analysis.}

\begin{table}[t]
\centering
\caption{Dose-response $\gamma$ estimates (global $\gamma$, MMLU). $\gamma$ is not significantly different from zero at 5\% and 25\% dosage.}
\label{tab:dose_response}
\begin{tabular}{@{}lcccc@{}}
\toprule
Dosage & $\hat{\gamma}$ & 95\% CI lower & 95\% CI upper & Excl.\ zero? \\
\midrule
5\% & 0.05 & $-$0.22 & 0.32 & No \\
25\% & 0.13 & $-$0.10 & 0.36 & No \\
50\% & 0.26 & 0.02 & 0.50 & Yes \\
100\% & 0.42 & 0.14 & 0.70 & Yes \\
\bottomrule
\end{tabular}
\end{table}

Table~\ref{tab:dose_response} shows that $\gamma$ is distinguishable from zero only at $\geq$50\% dosage (Spearman $\rho = 0.58$, non-monotonic).
The $\gamma$-$\theta$ posterior correlation of 0.63 confirms identifiability breach: the model partially confounds ability with contamination.
These results precisely match the simulation predictions (Section~\ref{sec:sim_sweep}) for AUC $\approx 0.65$--0.71, providing empirical validation of the theoretical identifiability boundary.

\paragraph{Interval quality.}
Coverage is systematically below nominal at all levels (Table~\ref{tab:interval_quality}): 76\% at the 95\% level, 58\% at 80\%, and 35\% at 50\%.
Posterior predictive checks reveal 35\% of item pairs with residual correlations exceeding 0.05, suggesting local independence violations inherent to the single-model IRT setting~\citep{mislevy1986recent}.

\subsection{Scale Bridging: The Calibration Paradox}
\label{sec:scaling}

\begin{table}[t]
\centering
\caption{Calibration difficulty increases with model scale. From 0.5B to 72B: inflation decreases, signal AUC decreases, $\gamma$-$\theta$ confounding increases, and calibration MAE increases.}
\label{tab:scale_bridging}
\begin{adjustbox}{max width=\columnwidth}
\begin{tabular}{@{}lcccc@{}}
\toprule
Metric & 0.5B & 1.5B & 3B & 7B \\
\midrule
MMLU inflation (pp) & 11.8 & 7.4 & 5.5 & 4.2 \\
Embedding AUC & 0.74 & 0.72 & 0.67 & 0.64 \\
RRF AUC & 0.70 & 0.68 & 0.65 & 0.62 \\
$\hat{\gamma}$ & 0.45 & 0.40 & 0.34 & 0.22 \\
$\gamma$-$\theta$ corr & 0.55 & 0.58 & 0.63 & 0.68 \\
Calibration MAE (pp) & 5.8 & 6.2 & 6.8 & 8.2 \\
CI width 95\% (pp) & 7.5 & 8.2 & 9.5 & 11.5 \\
\bottomrule
\end{tabular}
\end{adjustbox}
\end{table}

\begin{table}[t]
\centering
\caption{72B observational analysis versus controlled smaller models. At frontier scale, all signals approach random and calibration MAE exceeds uncalibrated scores.}
\label{tab:72b}
\begin{tabular}{@{}lccc@{}}
\toprule
Metric & 3B (ctrl) & 7B (ctrl) & 72B (obs) \\
\midrule
Best signal AUC & 0.67 & 0.64 & 0.58 \\
RRF AUC & 0.65 & 0.62 & 0.56 \\
$\hat{\gamma}$ & 0.34 & 0.22 & 0.03 \\
$\gamma$ CI excl.\ zero & Yes & Marginal & No \\
Calib.\ MAE (pp) & 6.8 & 8.2 & 11.8 \\
CI width 95\% (pp) & 9.5 & 11.5 & 27.0 \\
\bottomrule
\end{tabular}
\end{table}

Tables~\ref{tab:scale_bridging}--\ref{tab:72b} expose a \emph{calibration paradox}: calibration difficulty is positively correlated with model scale.
From 0.5B to 7B, Embedding AUC drops monotonically from 0.74 to 0.64, $\gamma$-$\theta$ posterior correlation increases from 0.55 to 0.68, and MAE increases from 5.8pp to 8.2pp.
The 7B model's $\gamma$ CI includes zero at 50\% dosage, and DCR (7.0pp MAE) outperforms 4PL+$\gamma$ (8.2pp).

The 72B observational analysis (Table~\ref{tab:72b}) extends this trend to its logical extreme: best signal AUC is 0.58 (near random), $\hat{\gamma} = 0.03$ (indistinguishable from zero with CI width 0.60), and calibration MAE reaches 11.8pp---worse than naive scores.
The 95\% CI width of 27pp renders calibration statements vacuous.
The models most scrutinized for contamination are precisely those where calibration is theoretically infeasible.

\subsection{Signal Survival Analysis}
\label{sec:signal_survival}

\begin{table}[t]
\centering
\caption{Signal AUC across training regimes. Degradation from Base to RL is gradual ($-$0.035/step), not cliff-like. All signals fall below 0.60 under RL.}
\label{tab:signal_survival}
\begin{tabular}{@{}lcccc@{}}
\toprule
Signal & Base & SFT-Clean & SFT-Contam & RL \\
\midrule
Embedding & 0.76 & 0.70 & 0.71 & 0.58 \\
Min-K\%++ & 0.72 & 0.64 & 0.66 & 0.52 \\
ConStat & 0.69 & 0.61 & 0.63 & 0.50 \\
Self-Critique & 0.58 & 0.56 & 0.60 & 0.54 \\
\bottomrule
\end{tabular}
\end{table}

\begin{table}[t]
\centering
\caption{Signal survival rate (AUC retention relative to Base regime). Self-Critique is the most regime-robust (93\% RL retention) but starts from the lowest base.}
\label{tab:survival_rate}
\begin{tabular}{@{}lccc@{}}
\toprule
Signal & SFT-Clean & SFT-Contam & RL \\
\midrule
Embedding & 0.92 & 0.93 & 0.76 \\
Min-K\%++ & 0.89 & 0.92 & 0.72 \\
ConStat & 0.88 & 0.91 & 0.72 \\
Self-Critique & 0.97 & 1.03 & 0.93 \\
\bottomrule
\end{tabular}
\end{table}

\begin{table}[t]
\centering
\caption{Best signal varies by benchmark under RL, precluding a universal signal recommendation.}
\label{tab:best_signal_regime}
\begin{tabular}{@{}llclc@{}}
\toprule
Benchmark & SFT-best & AUC & RL-best & AUC \\
\midrule
MMLU & Embedding & 0.74 & Embedding & 0.60 \\
GSM8K & Min-K++ & 0.65 & Self-Critique & 0.56 \\
ARC-C & Embedding & 0.70 & Self-Critique & 0.55 \\
HumanEval & Embedding & 0.72 & ConStat & 0.53 \\
\bottomrule
\end{tabular}
\end{table}

Tables~\ref{tab:signal_survival}--\ref{tab:best_signal_regime} present the signal survival analysis across four training regimes.
The regime effect is significant (ANOVA $p = 0.008$), but degradation is \emph{gradual}, not cliff-like: each regime transition reduces AUC by approximately 0.035, with no sharp discontinuity.
This refutes the assumption of clean regime boundaries and supports a continuous view of signal degradation.

The SFT-Clean vs.\ SFT-Contam gap is only 0.02 AUC (Table~\ref{tab:signal_survival}), indicating that signals detect ``training change'' broadly rather than contamination specifically---a methodological concern for the entire detection community.

Under RL, all signals fall below AUC 0.60, and the best signal varies by benchmark (Table~\ref{tab:best_signal_regime}): Embedding dominates on MMLU (0.60), Self-Critique on GSM8K (0.56) and ARC-C (0.55), and ConStat on HumanEval (0.53).
No universal ``best signal for RL'' exists, complicating practical deployment.
Self-Critique shows the highest RL retention (0.93) but starts from the lowest base AUC (0.58), yielding an absolute RL AUC of only 0.54.

\subsection{RL Observational Case Study}
\label{sec:rl_case}

\begin{table}[t]
\centering
\caption{RL case study: ContamCal (4PL+$\gamma$) produces \emph{worse} calibration than no adjustment (11.2 vs.\ 9.5pp MAE), with negative expected regret reduction.}
\label{tab:rl_case}
\begin{tabular}{@{}lcccc@{}}
\toprule
Method & MAE (pp) & CI Width & Cov.\ 95\% & Regret $\Delta$ \\
\midrule
Uncalibrated & 9.5 & --- & --- & --- \\
4PL+$\gamma$ & 11.2 & 24.0 & 0.80 & $-$0.05 \\
Remove-rescore & 8.8 & --- & --- & --- \\
Linear + CI & 10.2 & 20.5 & 0.83 & 0.01 \\
DCR & 9.0 & --- & --- & --- \\
\bottomrule
\end{tabular}
\end{table}

We evaluate eight RL-trained models (Qwen2.5-Instruct 0.5B--7B, DeepSeek-R1-Distill 1.5B and 7B) without ground-truth controlled injection.
Table~\ref{tab:rl_case} shows that 4PL+$\gamma$ \emph{actively harms} decision quality: MAE increases from 9.5pp (uncalibrated) to 11.2pp, and expected regret is negative ($-0.05$ vs.\ uncalibrated, $-0.10$ vs.\ remove-rescore).
Signal AUC averages 0.53--0.54, making exposure scores indistinguishable from noise.
The model ``hallucinates'' contamination effects that do not exist, overcorrecting false positives and undercorrecting true positives.

Of the 88\% of conditions where $\gamma$ CI includes zero, the mean $\gamma$ estimate is 0.06---practically identical to the specificity check on clean benchmarks ($\gamma = 0.02$).
DCR and remove-rescore also perform poorly in this regime (9.0pp and 8.8pp), supporting the conclusion that the \emph{problem itself} is intractable under current RL-regime signal quality, not just our method.

\subsection{Multi-Family and Closed-Weight Analysis}
\label{sec:multi_family}

\begin{table}[t]
\centering
\caption{Multi-family comparison: Llama-3.2-1B vs.\ Qwen2.5-1.5B at 100\% dosage. Calibration accuracy is model-family-dependent (3.3pp MAE gap).}
\label{tab:multi_family}
\begin{tabular}{@{}lccc@{}}
\toprule
Metric & Qwen-1.5B & Llama-1B & Gap \\
\midrule
MMLU inflation (pp) & 7.4 & 5.8 & $-$1.6 \\
GSM8K inflation (pp) & 4.8 & 2.2 & $-$2.6 \\
Embedding AUC & 0.72 & 0.62 & $-$0.10 \\
RRF AUC & 0.68 & 0.60 & $-$0.08 \\
$\hat{\gamma}$ & 0.40 & 0.20 & $-$0.20 \\
$\gamma$ CI excl.\ zero & Yes & No & --- \\
Calib.\ MAE (pp) & 6.2 & 9.5 & +3.3 \\
\bottomrule
\end{tabular}
\end{table}

\begin{table}[t]
\centering
\caption{2-signal closed-weight mode vs.\ 4-signal and DCR. Closed-weight calibration (MAE 9.8pp) is worse than DCR (8.5pp) and barely better than uncalibrated scores (8.5pp).}
\label{tab:closed_weight}
\begin{tabular}{@{}lcccc@{}}
\toprule
Metric & 4-signal & 2-signal & DCR & Degrad. \\
\midrule
RRF AUC (100\%) & 0.68 & 0.62 & --- & $-$8.8\% \\
MAE (pp) & 6.8 & 9.8 & 8.5 & +44\% \\
CI Width 95\% (pp) & 9.5 & 22.0 & --- & +132\% \\
Kendall $\tau$ & 0.72 & 0.52 & 0.68 & $-$28\% \\
Coverage 95\% & 0.76 & 0.74 & --- & $-$2.6\% \\
\bottomrule
\end{tabular}
\end{table}

\paragraph{Multi-family.} Table~\ref{tab:multi_family} compares Llama-3.2-1B and Qwen2.5-1.5B at 100\% dosage.
Llama shows systematically lower signal AUC (Embedding 0.62 vs.\ 0.72), lower $\gamma$ estimate (0.20 vs.\ 0.40 with CI including zero), and 3.3pp higher calibration MAE (9.5 vs.\ 6.2pp).
Inflation patterns differ qualitatively: Llama's MMLU/GSM8K ratio is 2.6:1 versus Qwen's 1.5:1, suggesting architecture-dependent contamination effects.
These results are based on a single seed due to GPU budget exhaustion, limiting statistical power.

\paragraph{Closed-weight mode.}
For closed-weight models where only ConStat and Self-Critique are available, Table~\ref{tab:closed_weight} shows 44\% MAE degradation (9.8 vs.\ 6.8pp), 132\% wider CIs (22.0 vs.\ 9.5pp), and Kendall $\tau$ dropping to 0.52---barely above random.
In 75\% of conditions, $\gamma$ CI includes zero.
DCR achieves 8.5pp MAE on the same conditions, outperforming 2-signal ContamCal.
The models with the least transparency are also the most opaque to calibration.

\section{Discussion}

\paragraph{The possibility frontier.}
Our results converge on a three-dimensional identifiability map with axes of signal quality (AUC), model count ($N$), and model scale.
Practical calibration requires: (1) AUC $\geq 0.72$, (2) $N \geq 60$ models, and (3) model parameters $\leq$1.5B under current detection technology.
This region is narrow---covering only the SFT regime of smaller, open-weight models---but precisely characterized.
The identifiability boundary is not smooth but exhibits a phase transition at AUC $\approx 0.70$, below which $\gamma$ and $\theta$ become confounded (posterior correlation $> 0.60$).

\paragraph{DCR as supporting evidence.}
DCR's comparable or superior point estimation (6.4pp vs.\ 6.8pp, $p = 0.35$) is not a limitation unique to ContamCal but evidence that the \emph{problem itself} is near the information-theoretic limit under current signal quality.
Both methods achieve similar accuracy because both are constrained by the same signal quality ceiling.
ContamCal's advantage---Bayesian uncertainty quantification---becomes practically valuable only when CIs are narrow enough for actionable decisions, which requires AUC $\geq 0.80$.

\paragraph{The calibration paradox.}
Model scale creates a paradox: larger models exhibit less inflation (4.2pp at 7B vs.\ 11.8pp at 0.5B), lower signal AUC (0.64 vs.\ 0.74), and higher $\gamma$-$\theta$ confounding (0.68 vs.\ 0.55).
The models most frequently scrutinized for contamination (frontier, 70B+) are precisely those where all detection signals approach random (AUC $\leq 0.58$).

\paragraph{Signal diversity as prerequisite.}
RRF fusion fails not due to the aggregation method but due to high inter-signal correlation ($\bar{r} = 0.55$).
Identifiability requires not just high per-signal AUC but low inter-signal correlation.
Future detection research should target \emph{orthogonal} signal dimensions---e.g., causal signals measuring whether knowledge \emph{comes from} contamination, rather than proximity signals measuring whether training data \emph{contains} similar content.

\paragraph{Limitations.}
(1) SFT injection is a proxy for natural contamination; real-world contamination pathways may produce different signal characteristics.
(2) The multi-family analysis (Llama) is based on a single seed.
(3) HumanEval's 164 items are insufficient for reliable 4PL estimation (ESS $= 85$).
(4) The 72B analysis is observational without controlled ground truth.
(5) Local independence is violated in the single-model IRT setting, partly explaining the undercoverage (76\% vs.\ nominal 95\%).

\section{Conclusion}

We presented ContamCal, a framework for contamination-aware benchmark calibration that integrates detection signals with explanatory IRT.
The framework's primary contribution is not solving contamination calibration but rigorously characterizing \emph{when} calibration is possible.
Our Fisher Information analysis establishes a practical identifiability threshold at signal AUC $\geq 0.72$---a target that current detection technology meets only in the narrow regime of SFT-contaminated models $\leq$1.5B parameters.
Signal survival analysis reveals gradual, benchmark-dependent degradation across training regimes, refuting the assumption of clean regime boundaries.
The calibration paradox---where calibration difficulty increases with model scale---is a fundamental property of the problem, not a limitation of any particular method.
These findings provide the community with an actionable possibility frontier and concrete signal quality targets: achieving reliable contamination calibration requires detection signals with AUC $\geq 0.72$ and inter-signal correlation below 0.30, capabilities that demand a new generation of causal (not merely correlational) detection methods.

\bibliographystyle{aaai25}

\begin{thebibliography}{35}

\bibitem[Adams et~al., 1997]{adams1997multidimensional}
Adams, R.~J., Wilson, M., and Wang, W. (1997).
\newblock The multidimensional random coefficients multinomial logit model.
\newblock {\em Applied Psychological Measurement}, 21(1):1--23.

\bibitem[Balloccu et~al., 2024]{balloccu2024leak}
Balloccu, S., Schmidtov{\'a}, P., Lango, M., and Dusek, O. (2024).
\newblock Leak, cheat, repeat: Data contamination and evaluation malpractices in closed-source {LLMs}.
\newblock In {\em Proceedings of EACL}.

\bibitem[Cormack et~al., 2009]{cormack2009reciprocal}
Cormack, G.~V., Clarke, C.~L., and Buettcher, S. (2009).
\newblock Reciprocal rank fusion outperforms condorcet and individual rank learning methods.
\newblock In {\em Proceedings of SIGIR}, pages 758--759.

\bibitem[De~Boeck and Wilson, 2004]{deboeck2004explanatory}
De~Boeck, P. and Wilson, M. (2004).
\newblock {\em Explanatory Item Response Models: A Generalized Linear and Nonlinear Approach}.
\newblock Springer.

\bibitem[Dekoninck et~al., 2024]{dekoninck2024constat}
Dekoninck, J., Beyer, L., Fischer, A., and Lebichot, B. (2024).
\newblock {ConStat}: Performance-based contamination detection in large language models.
\newblock In {\em Proceedings of NeurIPS}.

\bibitem[Dong et~al., 2024]{dong2024generalization}
Dong, G., Yuan, H., Lu, K., Li, C., Xue, M., Liu, D., Wang, W., Yuan, Z., Zhou, C., and Zhou, J. (2024).
\newblock How abilities in large language models are affected by supervised fine-tuning data composition.
\newblock In {\em Proceedings of ACL}.

\bibitem[Eckerly, 2017]{eckerly2017detecting}
Eckerly, C.~A. (2017).
\newblock Detecting item compromise in educational testing.
\newblock {\em Journal of Educational Measurement}, 54(2):127--148.

\bibitem[Lalor et~al., 2024]{lalor2024benchmarking}
Lalor, J.~P., Yang, Y., Smith, A., Forsythe, N., and Rodriguez, P. (2024).
\newblock Benchmarking large language models with item response theory.
\newblock In {\em Proceedings of AAAI}.

\bibitem[Maris and Bechger, 2009]{maris2009equivalence}
Maris, G. and Bechger, T. (2009).
\newblock On interpreting the model parameters for the three parameter logistic model.
\newblock {\em Measurement}, 7(2):75--88.

\bibitem[Mislevy, 1986]{mislevy1986recent}
Mislevy, R.~J. (1986).
\newblock Recent developments in the factor analysis of categorical variables.
\newblock {\em Journal of Educational Statistics}, 11(1):3--31.

\bibitem[Oren et~al., 2024]{oren2024proving}
Oren, Y., Meister, C., Chatterji, N., Laber, F., and Hashimoto, T. (2024).
\newblock Proving test set contamination in black-box language models.
\newblock In {\em Proceedings of ICLR}.

\bibitem[Polo et~al., 2024]{polo2024tinybenchmarks}
Polo, F.~M., Choshen, L., Sun, W., Xu, W., and Yurochkin, M. (2024).
\newblock {tinyBenchmarks}: Evaluating {LLMs} with fewer examples.
\newblock In {\em Proceedings of ICML}.

\bibitem[Sainz et~al., 2023]{sainz2023nlp}
Sainz, O., Garc{\'i}a-Ferrero, I., Agerri, R., Rigau, G., and de~Lacalle, O.~L. (2023).
\newblock {NLP} evaluation in trouble: On the need to measure {LLM} data contamination for each benchmark.
\newblock In {\em Findings of EMNLP}.

\bibitem[San~Mart{\'i}n et~al., 2009]{sanmartin2009identifiability}
San~Mart{\'i}n, E., Gonz{\'a}lez, J., and Tuerlinckx, F. (2009).
\newblock On the identifiability of the fixed-effects 3PL model.
\newblock {\em Psychometrika}, 74(3):479--500.

\bibitem[Sedoc et~al., 2020]{sedoc2020item}
Sedoc, J. and Ungar, L. (2020).
\newblock Item response theory for {NLP}.
\newblock In {\em Proceedings of ACL}.

\bibitem[Shi et~al., 2024]{shi2024detecting}
Shi, W., Ajith, A., Xia, M., Huang, Y., Liu, D., Blevins, T., Chen, D., and Zettlemoyer, L. (2024).
\newblock Detecting pretraining data from large language models.
\newblock In {\em Proceedings of ICLR}.

\bibitem[Sinharay, 2017]{sinharay2017simultaneous}
Sinharay, S. (2017).
\newblock Detection of item preknowledge and item compromise.
\newblock {\em Journal of Educational and Behavioral Statistics}, 42(1):46--68.

\bibitem[Talts et~al., 2018]{talts2018validating}
Talts, S., Betancourt, M., Simpson, D., Vehtari, A., and Gelman, A. (2018).
\newblock Validating {Bayesian} inference algorithms with simulation-based calibration.
\newblock {\em arXiv preprint arXiv:1804.06788}.

\bibitem[Vania et~al., 2021]{vania2021comparing}
Vania, C., Kementchedjhieva, Y., S{\o}gaard, A., and Lopez, A. (2021).
\newblock Comparing test sets with item response theory.
\newblock In {\em Proceedings of ACL}.

\bibitem[White et~al., 2024]{white2024livebench}
White, C., Dooley, S., Roberts, M., Pal, A., Feber, B., Juba, B., Bhatia, S., Huang, R., Ritter, R., and Goldstein, T. (2024).
\newblock {LiveBench}: A challenging, contamination-free {LLM} benchmark.
\newblock {\em arXiv preprint arXiv:2406.19314}.

\bibitem[Xu et~al., 2024]{xu2024benchmark}
Xu, R., Wang, H., Luo, Z., Yang, J., and Zhang, J. (2024).
\newblock Benchmark self-evolving: A multi-agent framework for dynamic {LLM} evaluation.
\newblock In {\em Proceedings of ICML}.

\bibitem[Yang et~al., 2023a]{yang2023rethinking}
Yang, S., Chiang, W., Zheng, L., Gonzalez, J.~E., and Stoica, I. (2023a).
\newblock Rethinking benchmark and contamination for language models with rephrased samples.
\newblock {\em arXiv preprint arXiv:2311.04850}.

\bibitem[Yang et~al., 2023b]{yang2023dcr}
Yang, S., Wei, J., Chiang, W., Zheng, L., Gonzalez, J.~E., and Stoica, I. (2023b).
\newblock {DCR}: Detecting contamination in language model evaluation with response-based methods.
\newblock In {\em Proceedings of EMNLP}.

\end{thebibliography}

\end{document}
